\subsection{Algebraic closure of a field}

DEFINITION 3. --- Let K be a field. By an algebraic closure of K we understand any extension of K which is algebraic and algebraically closed.

Examples. --- * 1) The field C of complex numbers is an algebraic closure of the field R of real numbers (Gen.~Top., VIII, p.~100) *

\begin{enumerate}
\def\labelenumi{\arabic{enumi})}
\setcounter{enumi}{1}
\tightlist
\item
  Let K be a field and \(\Omega\) an algebraically closed extension of K. If \(\bar{K}\) is the relative algebraic closure of K in \(\Omega\) , then by V, p.~20, Prop. 2, \(\bar{K}\) is an algebraic closure of K. * In particular the field of all algebraic numbers (V, p.~20, Ex. 2) is an algebraic closure of the field Q of rational numbers. *
\end{enumerate}

PROPOSITION 6. --- Let \(\Omega\) be an extension of a field \(K\) . For \(\Omega\) to be an algebraic closure of \(K\) it is necessary and sufficient that it should be algebraic and that each non-constant polynomial in \(K[X]\) should split in \(\Omega[X]\) into a product of factors of degree 1.

The condition is necessary by (AC). Conversely, suppose that \(\Omega\) is algebraic over K and every non-constant polynomial over K{[}X{]} is a product in \(\Omega[\mathbf{X}]\) of factors of degree 1. Let \(\Omega'\) be an algebraic extension of \(\Omega\) and let \(x \in \Omega'\) . Since \(x\) is algebraic over \(\Omega\) and \(\Omega\) is algebraic over K, \(x\) is algebraic over K (V, p.~19, Prop. 3). Let \(f\) be the minimal polynomial of \(x\) over K. By hypothesis the polynomial \(f \in \mathbf{K}[\mathbf{X}]\) splits in \(\Omega[\mathbf{X}]\) into a product of factors of degree 1, whence \(x \in \Omega\) . Thus we have \(\Omega' = \Omega\) and \(\Omega\) is algebraically closed because it satisfies ( \(\mathrm{AC}'''\) ).

Remark 1. --- If \(\Omega\) is algebraic over \(K\) and if every non-constant polynomial of \(K[X]\) has a root in \(\Omega\) then \(\Omega\) is an algebraic closure of \(K\) (V, p.~156, Ex. 20).

PROPOSITION 7. --- Let \(\Omega\) be an algebraic extension of a field \(\mathbf{K}\) .

\begin{enumerate}
\def\labelenumi{\alph{enumi})}
\item
  If \(\Omega\) is algebraically closed, then every algebraic extension of \(\mathbf{K}\) is isomorphic to a subextension of \(\Omega\) .
\item
  Conversely suppose that every algebraic extension of finite degree of K is isomorphic to a subextension of \(\Omega\) ; then \(\Omega\) is algebraically closed.
\end{enumerate}

Assertion a) follows from Th. 1 (V, p.~20). Now assume the hypotheses of b) and consider a non-constant polynomial \(f \in K[X]\) . Let E be a splitting field of f (V, p.~21, Prop. 4); since E is algebraic of finite degree over K (V, p.~18, Th. 2), we may suppose that E is a subextension of \(\Omega\) . Then the polynomial f is a product of polynomials of degree 1 in \(\Omega[X]\) and Prop. 6 shows that \(\Omega\) is algebraically closed.

We can now prove the existence and uniqueness (up to isomorphism) of the algebraic closure of a field.

THEOREM 2 (Steinitz). --- Let K be a field ; then there exists an algebraic closure of K. If \(\Omega\) and \(\Omega'\) are two algebraic closures of K, there exists a K-isomorphism of \(\Omega\) onto \(\Omega'\) .

By Prop. 6 an algebraic closure of K is nothing other than a splitting extension for the set of all non-constant polynomials in K{[}X{]}. Th. 2 therefore follows from V, p.~21, Prop. 4 and V, p.~22, Cor.

COROLLARY. --- Let K and K' be two fields, Ω an algebraic closure of K and Ω' an algebraic closure of K'. For every isomorphism u of K onto K' there exists an isomorphism v of Ω onto Ω' extending u.

It is enough to apply Th. 2 to the algebraic closures \(\Omega\) and \((\Omega', u)\) of \(\mathbf{K}\) .

Remarks. --- 2) In the notation of the preceding Corollary there exist in general K-automorphisms of \(\Omega\) distinct from the identity. Hence there is in general no uniqueness about the isomorphism \(v\) of \(\Omega\) onto \(\Omega'\) extending the isomorphisms \(u\) of K onto K'. For similar reasons there is in general more than one isomorphism of a splitting extension E onto a splitting extension E' for the same family \((f_i)_{i \in I}\) of polynomials. We recall that by contrast, for the perfect closure we have uniqueness (V, p.~5).

\begin{enumerate}
\def\labelenumi{\arabic{enumi})}
\setcounter{enumi}{2}
\tightlist
\item
  Let K be a field and \(\Omega\) an algebraic closure of K. Then the following construction may be given for a splitting extension for a family \((f_{i})_{i\in I}\) of non-constant polynomials in K{[}X{]}: let \(R_{i}\) be the set of roots of \(f_{i}\) in \(\Omega\) and let \(R=\bigcup_{i\in I}R_{i}\) . Then \(K(R)\) is the unique subextension of \(\Omega\) which is a splitting
\end{enumerate}

extension for \((f_i)_{i\in I}\) (V, p.~22, Prop. 5).

\begin{enumerate}
\def\labelenumi{\arabic{enumi})}
\setcounter{enumi}{3}
\tightlist
\item
  Let K be a finite field and \(\Omega\) an algebraic closure of K. Then \(\Omega\) is infinite (V, p.~20, Prop. 3); since every extension of finite degree of K is a finite field, \(\Omega\) is an algebraic extension of infinite degree of K.
\end{enumerate}

